../cictikz/src/cictikz/data/tex/cictikz_preamble.tex

% Absolute Seebeck coefficient of eight metals against
% temperature, digitized from the Wikimedia original (Nanite, CC0) that
% this figure replaces. The span is the lecture's point: from about
% +20 uV/K for tungsten down to -60 uV/K for palladium, so a couple that
% pairs a positive metal with a negative one doubles the voltage per
% kelvin.
%
% Generated by py/tikzplot.py. Edit the script in ex/ that
% produced it and re-run that, not this file.

\begin{tikzpicture}[thick]

\begin{axis}[
  width=10.5cm,
  height=8.0cm,
  scale only axis,
  grid=both,
  grid style={gray!22, very thin},
  tick align=outside,
  tick label style={font=\small},
  label style={font=\small},
  title style={font=\small, yshift=-1mm},
  legend cell align=left,
  legend style={font=\small, draw=gray!50, fill=white, fill opacity=0.85, text opacity=1},
  xlabel={Temperature [K]},
  ylabel={Seebeck coef.\ [$\mu$V/K]}
]
\addplot[blue, very thick, mark=none] coordinates {(100,1.19) (150,1.12) (200,1.29) (273,1.7) (300,1.83) (400,2.34) (500,2.83) (600,3.33) (700,3.83) (800,4.34) (900,4.85) (1000,5.37) (1100,5.89) (1200,6.41) (1300,6.91)};
\addplot[armygreen, very thick, mark=none] coordinates {(100,0.73) (150,0.85) (200,1.05) (273,1.38) (300,1.51) (400,2.08) (500,2.82) (600,3.72) (700,4.72) (800,5.77) (900,6.85) (1000,7.95) (1100,9.06) (1200,10.15)};
\addplot[red, very thick, mark=none] coordinates {(100,0.82) (150,1.02) (200,1.34) (273,1.79) (300,1.94) (400,2.46) (500,2.86) (600,3.18) (700,3.43) (800,3.63) (900,3.77) (1000,3.85) (1100,3.88) (1200,3.86) (1300,3.78)};
\addplot[cyan, very thick, mark=none] coordinates {(100,4.29) (150,1.32) (200,-1.27) (273,-4.45) (300,-5.28) (400,-7.83) (500,-9.89) (600,-11.66) (700,-13.24) (800,-14.81) (900,-16.32) (1000,-17.79) (1100,-19.22) (1200,-20.62) (1300,-21.97) (1400,-23.32) (1600,-25.97) (1800,-28.57) (2000,-31.23)};
\addplot[magenta, very thick, mark=none] coordinates {(100,2) (150,-1.63) (200,-4.85) (273,-9) (300,-9.99) (400,-13) (500,-16.03) (600,-19.06) (700,-22.09) (800,-25.12) (900,-28.15) (1000,-31.18) (1100,-34.21) (1200,-37.24) (1300,-40.27) (1400,-43.3) (1600,-49.36) (1800,-55.42) (2000,-61.48)};
\addplot[olive, very thick, mark=none] coordinates {(273,0.13) (300,1.07) (400,4.44) (500,7.53) (600,10.29) (700,12.66) (800,14.65) (900,16.28) (1000,17.57) (1100,18.53) (1200,19.18) (1300,19.53) (1400,19.6) (1600,18.97) (1800,17.41) (2000,15.05) (2200,12.01) (2400,8.39)};
\addplot[black, very thick, mark=none] coordinates {(273,4.71) (300,5.57) (400,8.52) (500,11.12) (600,13.27) (700,14.94) (800,16.13) (900,16.86) (1000,17.16) (1100,17.08) (1200,16.65) (1300,15.92) (1400,14.94) (1600,12.42) (1800,9.52) (2000,6.67) (2200,4.3) (2400,2.87)};
\addplot[gray, very thick, mark=none] coordinates {(7,-0.2) (20,-0.63) (40,-0.66) (60,-0.71) (90,-0.83) (120,-0.92) (160,-1.03) (200,-1.1) (233,-1.18)};
\node[anchor=west, blue, scale=0.85, align=left] at (axis cs:1330,6.91) {Cu};
\node[anchor=west, armygreen, scale=0.85, align=left] at (axis cs:1230,10.15) {Ag};
\node[anchor=west, red, scale=0.85, align=left] at (axis cs:1330,3.78) {Au};
\node[anchor=west, cyan, scale=0.85, align=left] at (axis cs:2030,-31.23) {Pt};
\node[anchor=west, magenta, scale=0.85, align=left] at (axis cs:2030,-61.48) {Pd};
\node[anchor=west, olive, scale=0.85, align=left] at (axis cs:2430,8.39) {W};
\node[anchor=west, black, scale=0.85, align=left] at (axis cs:2430,2.87) {Mo};
\node[anchor=west, gray, scale=0.85, align=left] at (axis cs:263,-1.18) {Pb};
\end{axis}

\end{tikzpicture}
\end{document}
